% !TEX root = ~/Code/TN_LOWE_RECONSTRUCTION/main.tex
\section{Event Simulation}\label{sec:SIMULATION}
In a LArTPC, charged particles produced in neutrino interactions ionise the argon atoms. An applied electric field causes the liberated electrons to drift toward segmented readout planes, where their arrival times and positions are recorded, allowing for a 3D--reconstruction of the interaction. In addition to ionisation signals, LAr also produces scintillation light, which can be used for triggering and timing (reconstructing the drift coordinate). For low-energy events, efficient charge collection, background mitigation, and light detection are crucial for a reliable reconstruction and event selection. The neutrino event simulation in DUNE follows the standard low--energy workflows of the MC framework LArSoft for DUNE~\cite{Snider2017wjd}.

\subsection{Geometries\label{sec:GEOMETRY}}
In this work, we utilized two geometries: the \ac{HD} workspace geometry (hd\_1x2x6) and the \ac{VD} workspace geometry (vd\_1x8x14\_3view\_30deg). The \ac{HD} workspace features a single central \ac{APA} and a drift length of \SI{3.6}{\meter}, designed to represent the central region of the actual \ac{HD} detector. In contrast, the \ac{VD} workspace geometry consists of a single drift volume that is approximately \SI{7}{\meter} long. Both geometries aim to model the full DUNE Far Detector (FD) modules by representing a subvolume that can be scaled to the entire detector.

The \ac{HD} geometry is used in low-energy neutrino studies, simulating two background scenarios: one with complete external background simulated in the central \ac{APA} (hd\_1x2x6\_lateralAPA) and another without background in the central \ac{APA} (hd\_1x2x6\_centralAPA). In both scenarios, we simulate the full set of radiological backgrounds, along with top, bottom, head, and tail external backgrounds.

For the \ac{VD} geometry, we consider two scenarios: the nominal scenario, which likely represents the upper half of the \ac{VD} module (vd\_1x8x14\_3view\_30deg), and a background scenario that assumes additional shielding (vd\_1x8x14\_3view\_30deg\_shielded). The latter scenario is designed to represent the lower half of the \ac{VD} module, which is expected to have a different background profile due to the additional shielding.

\subsection{Neutrino Production\label{sec:NEUTRINO_GENERATION}}
The first step in simulating neutrino interactions in DUNE relies on a comprehensive nuclear model of the possible processes. For low--energy events, this model is \textbf{MARLEY} (Model of Argon Reaction Low Energy Yields), an event generator developed to simulate the charged-current (CC) interactions of low-energy neutrinos, typically in the range of a few to tens of \unit{\mega\electronvolt} with argon nuclei. It is particularly suited for simulating interactions relevant to solar and supernova neutrinos in liquid argon detectors, such as the CC process:

\begin{equation}
    \nu_e + {}^{40}\mathrm{Ar} \rightarrow e^- + {}^{40}\mathrm{K}^*
\end{equation}

The excitation of the residual nucleus \( {}^{40}\mathrm{K}^* \) leads to the emission of de-excitation \(\gamma\)-rays, which can produce observable signals (blips) in a LArTPC. MARLEY incorporates gamma--decay schemes to describe these processes with physically motivated detail.

% \iffalse
MARLEY is integrated into LArSoft (the Liquid Argon Software framework) and can be used to produce Monte Carlo (MC) samples for detector studies and physics analyses. It simulates:

\begin{itemize}
    \item Primary interaction kinematics
    \item Residual nuclear excited states
    \item Gamma-ray cascades from nuclear de-excitation
\end{itemize}
% \fi

\begin{figure}[ht!]
    \centering
    \includegraphics[width=\linewidth]{Images/hd_1x2x6_marley_signal_True_Electron_Energy.png}
    \caption[width=1\linewidth]{Neutrino--electron energy correlation in a CC interaction as predicted by MARLEY.\label{fig:NEUTRINO_ELECTRON_CORRELATION}}
\end{figure}

These MC events are passed to detector simulation (e.g., GEANT4-based tools) for full propagation of particles through the LAr volume and digitisation. This enables detailed studies of energy deposition, signal formation, and reconstruction performance for low-energy neutrino interactions.

\iffalse
    Because low-energy events typically produce short tracks and sparse activity, MARLEY-based simulations are essential for:

    \begin{enumerate}
        \item Developing reconstruction algorithms tailored to low-energy signatures.
        \item Estimating detection efficiency for solar and supernova neutrinos.
        \item Characterising backgrounds from radiological and cosmogenic sources.
        \item Validating event selection and signal extraction techniques.
    \end{enumerate}
\fi

\begin{figure}[ht!]
    \centering
    \includegraphics[width=.49\linewidth]{Images/hd_1x2x6_TMarleySumK_CC_Mean_Fractional_Channel_Energy.png}
    \includegraphics[width=.49\linewidth]{Images/hd_1x2x6_CC_Kinematic_Threshold.png}
    \caption{MARLEY's computation of CC neutrino interaction. (left) Daughter fraction of the original kinetic energy. (right) Average binding energy.\label{fig:MARLEY_CC_INTERACTION}}
\end{figure}

In~\cref{fig:MARLEY_CC_INTERACTION}, the average energy transfer and binding energy of the CC interaction are displayed. One can observe how the gamma channel starts dominating the interaction up to \SI{\sim10}{\mega\electronvolt}, and above that threshold, the main electron takes over the fractional energy, accompanied by a continuous increase in average kinematic threshold caused by the production of additional secondary particles. The most important secondary production is the emission of neutrons, which happens above \SI{\sim15}{\mega\electronvolt}, and leads to a decrease in visible energy that cannot be recovered under realistic reconstruction scenarios (with background). All of these effects combine to a strong energy dependence of the deposited and visible energies, which should be taken into account in our signal calibration. Given that the true neutrnio energy is our calibration goal, the aim of the following approach is to approximate a bin--by--bin calibration using the number of hits (NHits) in the main clusters as a proxy. Therefore, the following electron~\ref{sec:ELECTRON_CALIBRATION} and neutrino calibration~\ref{sec:NEUTRINO_CALIBRATION} steps divide the sample into regions of equal NHit threshold.
\subsection{Daughter Particle Propagation}
Following the generation of the primary neutrino interaction—e.g., via the MARLEY generator for low-energy solar neutrino events—the resulting final-state particles are propagated through the liquid argon medium using the \texttt{Geant4} toolkit. \texttt{Geant4} models the detailed particle-matter interactions, including energy loss via ionisation, multiple scattering, bremsstrahlung, photon production, and hadronic interactions, according to physics lists tailored for low-energy processes.

As charged particles traverse the liquid argon, they lose energy primarily through ionization, producing free ionization electrons and excitation of argon atoms, which subsequently decay and emit prompt scintillation photons in the vacuum ultraviolet (VUV) range (128~nm). These processes are central to the signal generation in a LArTPC. For computational efficiency, the simulation does not explicitly track every ionisation electron or VUV photon. Instead, energy loss along particle trajectories is recorded as a set of energy depositions, or \emph{edeps}, which include position, time, energy, and associated particle identity. Similarly, the number of scintillation photons produced in each energy deposition step is recorded as a secondary output, commonly referred to as \emph{simphotons}. These \texttt{edeps} and \texttt{simphotons} are then passed as inputs to a subsequent stage of simulation, often called \texttt{detsim}.

\subsection{Detector Response Simulation}
In this stage, effective propagation models are applied to simulate the transport of ionisation electrons and scintillation light to the TPC (wires or strips) and PDS (photodetectors) readout elements, incorporating effects such as recombination, diffusion, electron lifetime, light attenuation, and detector geometry. The detector response step models how the ionisation electrons and scintillation photons produced in the active volume give rise to signals on the readout systems. This stage takes as input the energy depositions (\texttt{edeps}) and the associated scintillation photons (\texttt{simphotons}) from the particle propagation step, and applies parametrised or semi-analytic models to simulate the detector's response.
%This two-stage simulation framework balances physical accuracy and computational efficiency, allowing for robust modeling of low-energy neutrino interactions in large-scale LArTPC detectors.

\subsubsection{TPC Signal}
Ionisation electrons drift under an applied electric field toward the TPC wire planes. During this drift, the electrons experience diffusion (both longitudinal and transverse), attachment to electronegative impurities (parametrised by the electron lifetime), and charge loss due to recombination. These effects are modelled to modify the total collected charge and the temporal and spatial spread of the signal. The LArTPC typically uses three wire planes to detect the drifting charge:

\begin{enumerate}
    \item\textbf{Collection plane} responds with unipolar signals as the electrons are finally absorbed.
    \item\textbf{Induction planes} respond with bipolar signals as electrons pass by without being collected.
\end{enumerate}

The motion of charge induces time-dependent signals (waveforms) on the wires, which are digitised at high frequency. These raw waveforms are then processed using electronics response models and overlaid with simulated electronic noise based on detector-specific measurements. A noise model is applied to simulate realistic signal-to-noise conditions and electronics shaping. To improve charge extraction, the waveforms undergo a \textbf{deconvolution procedure}, where the electronics response and noise spectrum are inverted in frequency space to recover a cleaner signal representation. The deconvolved waveforms are integrated over time windows to form localised \emph{hits}, which are stored with time, position (inferred from drift), and charge information.

\subsubsection{PDS Signal}

Simultaneously, the scintillation photons are detected by the Photon Detection System (PDS), typically composed of photomultiplier tubes (PMTs), SiPMs, or wavelength-shifting light guides. The \texttt{simphotons} from Geant4 are converted into photo-electrons using detector-specific quantum efficiencies and timing models.

The PDS response is also simulated with realistic timing resolution, dark noise, and single photoelectron gain characteristics. Simulated waveforms are produced for each optical channel, overlaid with noise, and optionally deconvolved or filtered depending on the readout electronics. Like the TPC signals, these waveforms are processed to extract \emph{optical hits}—integrated photon pulses with associated time and charge information.

\subsubsection{Monte Carlo Output}

The output of the \texttt{detsim} stage is a set of digitised and processed hits in both the TPC and PDS systems. These hits serve as the basis for downstream reconstruction, including clustering, tracking, and calorimetric energy reconstruction (see \cref{sec:RECONSTRUCTION}).

\iffalse
    \subsection{Signal Reconstruction}\label{sec:NEUTRINO_SIMULATION}
    The latest advancements of DUNE's solar analyses start working from isolated TPC and PDS hits. Higher--order objects are available in the standard reconstruction chain, but the small imprint of these low--energy interactions requires to evaluate electron track with similar sizes to the TPC pitch or nuclear de--excitation activity ("blip--like" information) which gets easily lost in higher order track-- or shower--like reconstruction algorithms.

    \subsection{Background Model}\label{sec:BACKGROUND_MODEL}
    To ensure a high signal-to-background ratio in solar neutrino analyses, these backgrounds must be characterised (and ideally suppressed) through a combination of:

    \begin{itemize}
        \item Detector shielding and overburden to reduce external radiation and cosmogenic activation.
        \item Careful material selection and screening to minimise component-induced radioactivity.
        \item Cryogenic purification and use of UAr to lower intrinsic radioisotope concentrations.
        \item Event-level discrimination using timing, topology, and energy reconstruction.
    \end{itemize}

    A detailed background model informed by in situ measurements, simulations, and radiopurity assays is essential to interpret the observed data with the required precision for solar neutrino physics.
\fi

\subsection{Monte Carlo Analyser}\label{sec:MC_ANALYSER}
As mentioned in the previous sections, the low energy simulation steps follow the standard generation, propagation and detector simulation steps of the MC workflow. The reconstruction step for these studies includes a basic clustering scheme that applies both to the TPC and PDS hits generating the final TPC clusters and PDS flashes to be used in the analysis. For the specific purpose of the solar neutrino studies an analyser module (\href{https://github.com/DUNE/duneana/blob/develop/duneana/SolarNuAna/SolarNuAna_module.cc}{SolarNuAna\_module.cc}) has been written to combine these reconstruction objects with the truth information of the event. This module essentially combines all the information needed to perform the calibration and reconstruction of the signal candidates. The biggest caveat from the analysis point of view is that all clusters and flash information is reduced to a point--like data structure. This is, of course, a choice of design in the analyser and allows for a simpler handling of the data because all event candidates have the same dimension. I future works, this analyser should be substituted by the more general \href{https://github.com/DUNE/duneana/blob/develop/duneana/LowEAna/LowEAna_module.cc}{LowEAna\_module.cc} which aims to provide the same event reconstruction while keeping the internal structure of the candidate clusters to allow for better calibration and discrimination strategies.